\IfFileExists{stacks-project.cls}{%
\documentclass{stacks-project}
}{%
\documentclass{amsart}
}

% For dealing with references we use the comment environment
\usepackage{verbatim}
\newenvironment{reference}{\comment}{\endcomment}
%\newenvironment{reference}{}{}
\newenvironment{slogan}{\comment}{\endcomment}
\newenvironment{history}{\comment}{\endcomment}

% For commutative diagrams we use Xy-pic
\usepackage[all]{xy}

% We use 2cell for 2-commutative diagrams.
\xyoption{2cell}
\UseAllTwocells

% We use multicol for the list of chapters between chapters
\usepackage{multicol}

% This is generally recommended for better output
\usepackage{lmodern}
\usepackage[T1]{fontenc}

% For cross-file-references
\usepackage{xr-hyper}

% Package for hypertext links:
\usepackage{hyperref}

% For any local file, say "hello.tex" you want to link to please
% use \externaldocument[hello-]{hello}
\externaldocument[introduction-]{introduction}
\externaldocument[conventions-]{conventions}
\externaldocument[sets-]{sets}
\externaldocument[categories-]{categories}
\externaldocument[topology-]{topology}
\externaldocument[sheaves-]{sheaves}
\externaldocument[sites-]{sites}
\externaldocument[stacks-]{stacks}
\externaldocument[fields-]{fields}
\externaldocument[algebra-]{algebra}
\externaldocument[brauer-]{brauer}
\externaldocument[homology-]{homology}
\externaldocument[derived-]{derived}
\externaldocument[simplicial-]{simplicial}
\externaldocument[more-algebra-]{more-algebra}
\externaldocument[smoothing-]{smoothing}
\externaldocument[modules-]{modules}
\externaldocument[sites-modules-]{sites-modules}
\externaldocument[injectives-]{injectives}
\externaldocument[cohomology-]{cohomology}
\externaldocument[sites-cohomology-]{sites-cohomology}
\externaldocument[dga-]{dga}
\externaldocument[dpa-]{dpa}
\externaldocument[sdga-]{sdga}
\externaldocument[hypercovering-]{hypercovering}
\externaldocument[schemes-]{schemes}
\externaldocument[constructions-]{constructions}
\externaldocument[properties-]{properties}
\externaldocument[morphisms-]{morphisms}
\externaldocument[coherent-]{coherent}
\externaldocument[divisors-]{divisors}
\externaldocument[limits-]{limits}
\externaldocument[varieties-]{varieties}
\externaldocument[topologies-]{topologies}
\externaldocument[descent-]{descent}
\externaldocument[perfect-]{perfect}
\externaldocument[more-morphisms-]{more-morphisms}
\externaldocument[flat-]{flat}
\externaldocument[groupoids-]{groupoids}
\externaldocument[more-groupoids-]{more-groupoids}
\externaldocument[etale-]{etale}
\externaldocument[chow-]{chow}
\externaldocument[intersection-]{intersection}
\externaldocument[pic-]{pic}
\externaldocument[weil-]{weil}
\externaldocument[adequate-]{adequate}
\externaldocument[dualizing-]{dualizing}
\externaldocument[duality-]{duality}
\externaldocument[discriminant-]{discriminant}
\externaldocument[derham-]{derham}
\externaldocument[local-cohomology-]{local-cohomology}
\externaldocument[algebraization-]{algebraization}
\externaldocument[curves-]{curves}
\externaldocument[resolve-]{resolve}
\externaldocument[models-]{models}
\externaldocument[functors-]{functors}
\externaldocument[equiv-]{equiv}
\externaldocument[pione-]{pione}
\externaldocument[etale-cohomology-]{etale-cohomology}
\externaldocument[proetale-]{proetale}
\externaldocument[relative-cycles-]{relative-cycles}
\externaldocument[more-etale-]{more-etale}
\externaldocument[trace-]{trace}
\externaldocument[crystalline-]{crystalline}
\externaldocument[spaces-]{spaces}
\externaldocument[spaces-properties-]{spaces-properties}
\externaldocument[spaces-morphisms-]{spaces-morphisms}
\externaldocument[decent-spaces-]{decent-spaces}
\externaldocument[spaces-cohomology-]{spaces-cohomology}
\externaldocument[spaces-limits-]{spaces-limits}
\externaldocument[spaces-divisors-]{spaces-divisors}
\externaldocument[spaces-over-fields-]{spaces-over-fields}
\externaldocument[spaces-topologies-]{spaces-topologies}
\externaldocument[spaces-descent-]{spaces-descent}
\externaldocument[spaces-perfect-]{spaces-perfect}
\externaldocument[spaces-more-morphisms-]{spaces-more-morphisms}
\externaldocument[spaces-flat-]{spaces-flat}
\externaldocument[spaces-groupoids-]{spaces-groupoids}
\externaldocument[spaces-more-groupoids-]{spaces-more-groupoids}
\externaldocument[bootstrap-]{bootstrap}
\externaldocument[spaces-pushouts-]{spaces-pushouts}
\externaldocument[spaces-chow-]{spaces-chow}
\externaldocument[groupoids-quotients-]{groupoids-quotients}
\externaldocument[spaces-more-cohomology-]{spaces-more-cohomology}
\externaldocument[spaces-simplicial-]{spaces-simplicial}
\externaldocument[spaces-duality-]{spaces-duality}
\externaldocument[formal-spaces-]{formal-spaces}
\externaldocument[restricted-]{restricted}
\externaldocument[spaces-resolve-]{spaces-resolve}
\externaldocument[formal-defos-]{formal-defos}
\externaldocument[defos-]{defos}
\externaldocument[cotangent-]{cotangent}
\externaldocument[examples-defos-]{examples-defos}
\externaldocument[algebraic-]{algebraic}
\externaldocument[examples-stacks-]{examples-stacks}
\externaldocument[stacks-sheaves-]{stacks-sheaves}
\externaldocument[criteria-]{criteria}
\externaldocument[artin-]{artin}
\externaldocument[quot-]{quot}
\externaldocument[stacks-properties-]{stacks-properties}
\externaldocument[stacks-morphisms-]{stacks-morphisms}
\externaldocument[stacks-limits-]{stacks-limits}
\externaldocument[stacks-cohomology-]{stacks-cohomology}
\externaldocument[stacks-perfect-]{stacks-perfect}
\externaldocument[stacks-introduction-]{stacks-introduction}
\externaldocument[stacks-more-morphisms-]{stacks-more-morphisms}
\externaldocument[stacks-geometry-]{stacks-geometry}
\externaldocument[moduli-]{moduli}
\externaldocument[moduli-curves-]{moduli-curves}
\externaldocument[examples-]{examples}
\externaldocument[exercises-]{exercises}
\externaldocument[guide-]{guide}
\externaldocument[desirables-]{desirables}
\externaldocument[coding-]{coding}
\externaldocument[obsolete-]{obsolete}
\externaldocument[fdl-]{fdl}
\externaldocument[index-]{index}

% Theorem environments.
%
\theoremstyle{plain}
\newtheorem{theorem}[subsection]{Theorem}
\newtheorem{proposition}[subsection]{Proposition}
\newtheorem{lemma}[subsection]{Lemma}

\theoremstyle{definition}
\newtheorem{definition}[subsection]{Definition}
\newtheorem{example}[subsection]{Example}
\newtheorem{exercise}[subsection]{Exercise}
\newtheorem{situation}[subsection]{Situation}

\theoremstyle{remark}
\newtheorem{remark}[subsection]{Remark}
\newtheorem{remarks}[subsection]{Remarks}

\numberwithin{equation}{subsection}

% Macros
%
\def\lim{\mathop{\mathrm{lim}}\nolimits}
\def\colim{\mathop{\mathrm{colim}}\nolimits}
\def\Spec{\mathop{\mathrm{Spec}}}
\def\Hom{\mathop{\mathrm{Hom}}\nolimits}
\def\Ext{\mathop{\mathrm{Ext}}\nolimits}
\def\SheafHom{\mathop{\mathcal{H}\!\mathit{om}}\nolimits}
\def\SheafExt{\mathop{\mathcal{E}\!\mathit{xt}}\nolimits}
\def\Sch{\mathit{Sch}}
\def\Mor{\mathop{\mathrm{Mor}}\nolimits}
\def\Ob{\mathop{\mathrm{Ob}}\nolimits}
\def\Sh{\mathop{\mathit{Sh}}\nolimits}
\def\NL{\mathop{N\!L}\nolimits}
\def\CH{\mathop{\mathrm{CH}}\nolimits}
\def\proetale{{pro\text{-}\acute{e}tale}}
\def\etale{{\acute{e}tale}}
\def\QCoh{\mathit{QCoh}}
\def\Ker{\mathop{\mathrm{Ker}}}
\def\Im{\mathop{\mathrm{Im}}}
\def\Coker{\mathop{\mathrm{Coker}}}
\def\Coim{\mathop{\mathrm{Coim}}}

% Boxtimes
%
\DeclareMathSymbol{\boxtimes}{\mathbin}{AMSa}{"02}

%
% Macros for moduli stacks/spaces
%
\def\QCohstack{\mathcal{QC}\!\mathit{oh}}
\def\Cohstack{\mathcal{C}\!\mathit{oh}}
\def\Spacesstack{\mathcal{S}\!\mathit{paces}}
\def\Quotfunctor{\mathrm{Quot}}
\def\Hilbfunctor{\mathrm{Hilb}}
\def\Curvesstack{\mathcal{C}\!\mathit{urves}}
\def\Polarizedstack{\mathcal{P}\!\mathit{olarized}}
\def\Complexesstack{\mathcal{C}\!\mathit{omplexes}}
% \Pic is the operator that assigns to X its picard group, usage \Pic(X)
% \Picardstack_{X/B} denotes the Picard stack of X over B
% \Picardfunctor_{X/B} denotes the Picard functor of X over B
\def\Pic{\mathop{\mathrm{Pic}}\nolimits}
\def\Picardstack{\mathcal{P}\!\mathit{ic}}
\def\Picardfunctor{\mathrm{Pic}}
\def\Deformationcategory{\mathcal{D}\!\mathit{ef}}
\title{The self map given by a section is homotopic to the identity on the associated simplicial object}
\author{The Stacks Project Authors}
\date{stacks-project:simplicial.tex:local-snapshot-20261012; GNU FDL 1.2 or later}
\begin{document}
\maketitle
\noindent\textit{Editorial note.} The counted statement is the author's own lemma environment labelled lemma-fibre-products-simplicial-object-w-section (source lines 4873-4884); the proof environment begins on the second line after it and is its own proof, closing at source line 4906.\par\medskip
\noindent\textit{Editorial note.} The proof explicitly leaves the verification of the relations to the reader while pointing at the relations-homotopy lemma; the retained card reproduces the author's text as written and does not supply the omitted verification, which is recorded in the pairing note rather than silently completed.\par\medskip
\noindent\textit{Editorial note.} The retained preamble is the chapter's own minimal preamble (sources/stacks-simplicial/preamble.tex, taken with its real line count), a verbatim excerpt of the native Xy-pic setup, the native theorem declarations with the mathematical macros, and the project's own no-op reference environments; no macro was invented or redefined.\par\medskip
\noindent\textit{Editorial scope.} One original result of the Stacks Project's Simplicial Methods chapter, counted once: when f has a section, the degreewise self map (s f)\textasciicircum{}\{n+1\} on the fibre-product simplicial object is homotopic to the identity, so the object is homotopy equivalent to the constant one.\par\medskip
\noindent\textit{Allowed standard prerequisites (editorial).}\begin{itemize}
\item The fibre-product simplicial object. The simplicial object U with U\_n the (n+1)-fold fibre product of Y over X, built naively from f : Y to X when C has fibre products (the author's example-fibre-products-simplicial-object, retained as context).. It is the object of the statement (source lines 4877-4880).
\item Sections and the maps s f. A section s of f with f s = id, hence an endomorphism s f of Y whose iterates appear in each degree.. It produces the self map of the statement and the g\textasciicircum{}1 used in the homotopy (source lines 4876, 4887).
\item The relations characterising a homotopy of simplicial morphisms. A family h\_\{n,i\} defines a homotopy if and only if it satisfies the five relations with the face and degeneracy maps (the author's lemma-relations-homotopy and remark-homotopy-better, retained as context).. It is the criterion the written-down family is claimed to satisfy (source lines 4901-4903).
\item The functor of points. Describing a morphism of iterated fibre products by its effect on generalised points.. It is the language in which the homotopy is first described (source lines 4888-4897).
\end{itemize}\medskip
\begin{example}
\label{example-fibre-products-simplicial-object}
Suppose that $Y\to X$ is a morphism of $\mathcal{C}$ such that all
the fibred products $Y \times_X Y \times_X \ldots \times_X Y$ exist.
Then we set $U_n$ equal to the $(n + 1)$-fold fibre product,
and we let $\varphi : [n] \to [m]$ correspond to the map
(on ``coordinates'')
$(y_0, \ldots, y_m) \mapsto (y_{\varphi(0)}, \ldots, y_{\varphi(n)})$.
In other words, the map $U_0 = Y \to U_1 = Y \times_X Y$ is the
diagonal map. The two maps $U_1 = Y \times_X Y \to U_0 = Y$ are the
projection maps.
\end{example}

\begin{lemma}
\label{lemma-relations-homotopy}
In the situation above, we have the following relations:
\begin{enumerate}
\item We have $h_{n, 0} = b_n$ and $h_{n, n + 1} = a_n$.
\item We have $d^n_j \circ h_{n, i} = h_{n - 1, i - 1} \circ d^n_j$
for $i > j$.
\item We have $d^n_j \circ h_{n, i} = h_{n - 1, i} \circ d^n_j$
for $i \leq j$.
\item We have $s^n_j \circ h_{n, i} = h_{n + 1, i + 1} \circ s^n_j$
for $i > j$.
\item We have $s^n_j \circ h_{n, i} = h_{n + 1, i} \circ s^n_j$
for $i \leq j$.
\end{enumerate}
Conversely, given a system of maps $h_{n, i}$ satisfying the
properties listed above, then these define a morphism
$h$ which is a homotopy from $a$ to $b$.
\end{lemma}

\begin{proof}
Omitted. You can prove the last statement using the fact,
see Lemma \ref{lemma-face-degeneracy-category} that
to give a morphism of simplicial objects is the
same as giving a sequence of morphisms $h_n$ commuting
with all $d^n_j$ and $s^n_j$.
\end{proof}

\begin{remark}
\label{remark-homotopy-better}
Let $\mathcal{C}$ be any category (no assumptions whatsoever). Let
$U$ and $V$ be simplicial objects of $\mathcal{C}$. Let $a, b : U \to V$
be morphisms of simplicial objects of $\mathcal{C}$. A
{\it homotopy from $a$ to $b$} is given by
morphisms\footnote{In the literature, often the maps
$h_{n + 1, i} \circ s_i : U_n \to V_{n + 1}$ are used instead
of the maps $h_{n, i}$. Of course the relations these maps satisfy
are different from the ones in Lemma \ref{lemma-relations-homotopy}.}
$h_{n, i} : U_n \to V_n$, for $n \geq 0$, $i = 0, \ldots, n + 1$
satisfying the relations of Lemma \ref{lemma-relations-homotopy}.
As in Definition \ref{definition-homotopy} we say the morphisms $a$ and $b$
are {\it homotopic} if there exists a sequence of morphisms
$a = a_0, a_1, \ldots, a_n = b$ from $U$ to $V$ such that for each
$i = 1, \ldots, n$ there either exists a homotopy from $a_{i - 1}$ to $a_i$
or there exists a homotopy from $a_i$ to $a_{i - 1}$.
Clearly, if $F : \mathcal{C} \to \mathcal{C}'$ is any functor
and $\{h_{n, i}\}$ is a homotopy from $a$ to $b$, then
$\{F(h_{n, i})\}$ is a homotopy from $F(a)$ to $F(b)$.
Similarly, if $a$ and $b$ are homotopic, then $F(a)$ and $F(b)$
are homotopic.
Since the lemma says that the newer notion is the same
as the old one in case finite coproduct exist, we deduce
in particular that functors preserve the original notion
whenever both categories have finite coproducts.
\end{remark}

\begin{lemma}
\label{lemma-face-degeneracy-category}
The category $\Delta$ is the universal category
with objects $[n]$, $n \geq 0$ and morphisms
$\delta^n_j$ and $\sigma^n_j$ such that (a) every morphism is
a composition of these morphisms, (b) the relations
listed in Lemma \ref{lemma-relations-face-degeneracy} are satisfied,
and (c) any relation among the morphisms is a consequence of
those relations.
\end{lemma}

\begin{proof}
Omitted.
\end{proof}

\begin{definition}
\label{definition-homotopy}
Let $\mathcal{C}$ be a category having finite coproducts.
Suppose that $U$ and $V$ are two simplicial objects of $\mathcal{C}$.
Let $a, b : U \to V$ be two morphisms.
\begin{enumerate}
\item We say a morphism
$$
h : U \times \Delta[1] \longrightarrow V
$$
is a {\it homotopy from $a$ to $b$} if $a = h \circ e_0$ and
$b = h \circ e_1$.
\item We say the morphisms $a$ and $b$ are {\it homotopic} or are
{\it in the same homotopy class}
if there exists a sequence of morphisms $a = a_0, a_1, \ldots, a_n = b$
from $U$ to $V$ such that for each $i = 1, \ldots, n$ there either exists
a homotopy from $a_{i - 1}$ to $a_i$ or there exists a homotopy
from $a_i$ to $a_{i - 1}$.
\end{enumerate}
\end{definition}

\begin{lemma}
\label{lemma-relations-face-degeneracy}
The morphisms $\delta^n_j$ and $\sigma^n_j$ satisfy the following relations.
\begin{enumerate}
\item If $0 \leq i < j \leq n + 1$, then
$\delta^{n + 1}_j \circ \delta^n_i =
\delta^{n + 1}_i \circ \delta^n_{j - 1}$.
In other words the diagram
$$
\xymatrix{
& [n] \ar[rd]^{\delta^{n + 1}_j} & \\
[n - 1] \ar[ru]^{\delta^n_i} \ar[rd]_{\delta^n_{j - 1}} & &
[n + 1] \\
& [n] \ar[ru]_{\delta^{n + 1}_i} &
}
$$
commutes.
\item If $0 \leq i < j \leq n - 1$, then
$\sigma^{n - 1}_j \circ \delta^n_i =
\delta^{n - 1}_i \circ \sigma^{n - 2}_{j - 1}$.
In other words the diagram
$$
\xymatrix{
& [n] \ar[rd]^{\sigma^{n - 1}_j} & \\
[n - 1] \ar[ru]^{\delta^n_i} \ar[rd]_{\sigma^{n - 2}_{j - 1}} & &
[n - 1] \\
& [n - 2] \ar[ru]_{\delta^{n - 1}_i} &
}
$$
commutes.
\item If $0 \leq j \leq n - 1$, then
$\sigma^{n - 1}_j \circ \delta^n_j = \text{id}_{[n - 1]}$
and
$\sigma^{n - 1}_j \circ \delta^n_{j + 1} = \text{id}_{[n - 1]}$.
In other words the diagram
$$
\xymatrix{
& [n] \ar[rd]^{\sigma^{n - 1}_j} & \\
[n - 1]
\ar[ru]^{\delta^n_j}
\ar[rd]_{\delta^n_{j + 1}}
\ar[rr]^{\text{id}_{[n - 1]}} & & [n - 1] \\
& [n] \ar[ru]_{\sigma^{n - 1}_j} &
}
$$
commutes.
\item If $0 < j + 1 < i \leq n$, then
$\sigma^{n - 1}_j \circ \delta^n_i =
\delta^{n - 1}_{i - 1} \circ \sigma^{n - 2}_j$.
In other words the diagram
$$
\xymatrix{
& [n] \ar[rd]^{\sigma^{n - 1}_j} & \\
[n - 1] \ar[ru]^{\delta^n_i} \ar[rd]_{\sigma^{n - 2}_j} & &
[n - 1] \\
& [n - 2] \ar[ru]_{\delta^{n - 1}_{i - 1}} &
}
$$
commutes.
\item If $0 \leq i \leq j \leq n - 1$, then
$\sigma^{n - 1}_j \circ \sigma^n_i =
\sigma^{n - 1}_i \circ \sigma^n_{j + 1}$.
In other words the diagram
$$
\xymatrix{
& [n] \ar[rd]^{\sigma^{n - 1}_j} & \\
[n + 1] \ar[ru]^{\sigma^n_i} \ar[rd]_{\sigma^n_{j + 1}} & &
[n - 1] \\
& [n] \ar[ru]_{\sigma^{n - 1}_i} &
}
$$
commutes.
\end{enumerate}
\end{lemma}

\begin{proof}
Omitted.
\end{proof}

\begin{lemma}
\label{lemma-fibre-products-simplicial-object-w-section}
Let $f : Y \to X$ be a morphism of a category $\mathcal{C}$ with
fibre products. Assume $f$ has a section $s$. Consider the
simplicial object $U$ constructed in
Example \ref{example-fibre-products-simplicial-object}
starting with $f$. The morphism $U \to U$ which in each degree
is the self map $(s \circ f)^{n + 1}$ of $Y \times_X \ldots \times_X Y$
given by $s \circ f$ on each factor is homotopic to the identity on $U$.
In particular, $U$ is homotopy equivalent to the constant
simplicial object $X$.
\end{lemma}

\begin{proof}
Set $g^0 = \text{id}_Y$ and $g^1 = s \circ f$.
We use the morphisms
\begin{eqnarray*}
Y \times_X \ldots \times_X Y \times \Mor([n], [1])
& \to &
Y \times_X \ldots \times_X Y \\
(y_0, \ldots, y_n) \times \alpha
& \mapsto &
(g^{\alpha(0)}(y_0), \ldots, g^{\alpha(n)}(y_n))
\end{eqnarray*}
where we use the functor of points point of view to define the maps.
Another way to say this is to say that
$h_{n, 0} = \text{id}$, $h_{n, n + 1} = (s \circ f)^{n + 1}$ and
$h_{n, i} = \text{id}_Y^{i + 1} \times (s \circ f)^{n + 1 - i}$.
We leave it to the reader to show that these satisfy the relations
of Lemma \ref{lemma-relations-homotopy}. Hence they define
the desired homotopy. See also Remark \ref{remark-homotopy-better}
which shows that
we do not need to assume anything else on the category $\mathcal{C}$.
\end{proof}

\end{document}
